\section{Experimental Setup}
\label{sec:experimental_setup}

Our evaluation follows a single shared training and evaluation
pipeline across all baselines and our method; the only difference
between configurations is the quantizer module used in the
autoencoder. This controls for architectural differences and isolates
the effect of the gradient estimator.

\subsection{Datasets}

Our main benchmark is CIFAR-10, a widely used dataset of 50,000
training and 10,000 validation natural images at $32 \times 32$
resolution across 10 classes. CIFAR-10 is small enough
to train multiple strong baselines to convergence in minutes on a
single H100 and large enough to exhibit the codebook-collapse regime
reliably; prior work in vector quantization reports on it as a
primary reconstruction benchmark~\citep{mentzer2024fsq}. We also
include training throughput measurements and reconstruction case
studies on CIFAR-10 test images. Scaling experiments on
higher-resolution datasets (FFHQ $256{\times}256$, ImageNet-1k
$128{\times}128$) are deferred to the appendix and to future work.

\subsection{Architecture}

We use a VQ-GAN-style encoder-decoder~\citep{esser2021vqgan,
vandenoord2017vqvae} simplified for the CIFAR-10 resolution. The
encoder applies a $3{\times}3$ convolution to $64$ channels,
followed by two residual blocks at each of three resolution levels
$\{32, 16, 8\}$ with channel multipliers $(1, 2, 2)$ and mid-block
self-attention at the bottleneck. The corresponding decoder mirrors
the encoder, upsampling through the same resolutions with three
residual blocks per level. Group-normalisation and SiLU activations
are used throughout. A $1{\times}1$ convolution projects the
encoder's 128-channel bottleneck to the quantizer's working
dimension $d = 8$, and a symmetric post-quantisation
$1{\times}1$ convolution lifts back to 128 channels before the
decoder. The resulting autoencoder has $5.40$M parameters. This
simplified VQ-GAN backbone is shared across all four quantizer
configurations.

\subsection{Quantizers}

We compare four quantizer modules that plug into the same encoder
and decoder. \textbf{VQ-STE} is the baseline of
\citet{vandenoord2017vqvae}, using the estimator
in~\eqref{eq:ste}. \textbf{\ours{}} is our proposed estimator
in~\eqref{eq:our-estimator}. \textbf{Gumbel-Softmax
VQ}~\citep{jang2017gumbel} replaces the hard $\arg\min$ with a
differentiable softmax over codes with temperature $\tau$ annealed
from $1.0$ to $0.5$ during training at rate $3{\cdot}10^{-5}$; a
$1{\times}1$ convolution maps the encoder bottleneck to the code
logits. \textbf{FSQ}~\citep{mentzer2024fsq} projects the bottleneck to
$d = 6$ channels, applies a tanh, and rounds each channel
independently to one of per-channel levels $(8,8,8,5,5,5)$, giving an
implicit codebook of $8^3 \cdot 5^3 = 64{,}000$ entries. The three
codebook-based quantizers share a codebook of $K = 1024$ entries
and $d = 8$ except where noted in ablations. VQ-STE and \ours{} use
the same EMA codebook update (decay $\gamma = 0.99$) with
dead-code reinitialisation from the current batch; this is a
standard recipe~\citep{razavi2019vqvae2} and is held identical across
both to ensure a fair comparison of the gradient estimator alone.

\subsection{Training}

All models are trained on a single NVIDIA H100 80GB GPU with PyTorch
2.6 and TF32 tensor cores. We use AdamW~\citep{loshchilov2019adamw}
with $\beta = (0.9, 0.99)$, weight decay $0$, gradient clipping at
$1.0$, and a cosine learning-rate schedule from the peak value to
$10^{-6}$. The batch size is $1024$ with a peak learning rate of
$1.6 \times 10^{-3}$ (linearly scaled from the $2 \times 10^{-4}$
baseline used with batch $128$). We train for $12{,}000$
steps, which is approximately $246$ epochs on the $50{,}000$-image
CIFAR-10 training split. The reconstruction loss is an $\ell_1$
error on pixel values in $[0, 1]$ (we do not add an LPIPS or
adversarial loss to avoid confounding the gradient-estimator
comparison with loss engineering). The commitment coefficient is
$\beta = 0.25$ for the codebook-based quantizers, matching
\citet{vandenoord2017vqvae}; FSQ has no explicit commitment loss.

\subsection{Evaluation metrics}

\paragraph{Reconstruction quality.}
We report peak signal-to-noise ratio (PSNR) and structural
similarity (SSIM) on the full $10{,}000$-image CIFAR-10 validation
set. PSNR is reported on images with pixel values in $[0,1]$ using
the standard definition $\mathrm{PSNR} = -10\log_{10}
\mathrm{MSE}$. SSIM is computed per-image via scikit-image and
averaged. For a subset of configurations we additionally report
LPIPS~\citep{zhang2018lpips} with the AlexNet backbone as a
perceptual complement.

\paragraph{Codebook statistics.}
For codebook-based quantizers we report the \emph{usage} fraction
(the number of codes used at least once on the validation set
divided by $K$), the \emph{perplexity}
$\exp(-\sum_k p_k \log p_k)$ where $p_k$ is the empirical probability
of code $k$ on the validation set, and the number of \emph{active
codes per batch} during training. A higher perplexity relative to
$K$ indicates a more uniform use of the codebook, which is a widely
accepted proxy for the ``capacity'' actually used by the
tokenizer~\citep{huh2023straighteningste,
zheng2023onlineclustered}. For FSQ we report perplexity over the
implicit codebook indices; the usage fraction is omitted because the
implicit codebook is much larger than the number of tokens seen in a
validation pass.

\paragraph{Throughput.}
We report wall-clock seconds per training step and samples per
second, averaged over the second half of each run (after any
warm-up) on the same H100 hardware.

\subsection{Ablations}

In addition to the main comparison we run two ablation axes within
\ours{}. The first is the \emph{gradient-routing ablation} over
\{\textsc{ste}, rotation-only, rescale-only, full\} of
Section~\ref{sec:method}, which isolates the contribution of the
rotation $\rmat$ and the rescaling $s$ while keeping all other
training hyperparameters fixed. The second is the \emph{codebook-size
ablation} over $K \in \{1024, 4096, 8192, 16384\}$, which probes
whether the benefit of the rotation estimator grows or shrinks with
codebook capacity. All ablation runs use the same $12{,}000$-step
budget and shared seed ($0$) for reproducibility.

\subsection{Reproducibility}

We fix the PyTorch, NumPy, and CUDA random seeds at the start of each
run. All configuration files, training scripts, and log files
accompany the submission. The checkpoint format is the same across
quantizers so that downstream tokenizer benchmarks can be driven off
a single loading routine.
